% !TeX root = ../build/main.tex

An election involves four types of participants: organizer, \keywardens, voters, and sequencers. In addition, any party can follow an election as an observer.

\paragraph{Organizer.} The organizer is the entity responsible for defining and setting up the election. Its key responsibilities include defining the voting parameters, gathering the census data, and choosing who holds the election key. The organizer is identified by the Ethereum address that creates the election, and only that address can later change the parameters that the protocol allows to be changed (the status, the end time, the census of a dynamic census, the metadata, the maximum number of voters, and the grace window). The organizer ensures that the election is structured correctly but does not participate in vote collection or tallying.

\paragraph{\Keywardens.} \Keywardens are a set of decentralized parties that collaboratively generate a public encryption key that voters use to encrypt their votes. The corresponding secret key is secret-shared among them so that a threshold of key wardens can later collaborate to decrypt the final tally. Individual ballots remain private because the protocol only invokes threshold decryption on the final aggregate, not on each ballot separately. \Davinci supports two kinds of \keywardens: a \emph{\ac{DKG} committee}, drawn by lottery from a permissionless set of registered operators, which prepares a pool of keys in advance and assigns one to each election; and a \emph{council}, a small group of members invited by an organizer, which runs a dedicated key ceremony. When the organizer accepts the corresponding trust, a single sequencer can play the role of the \keywardens and hold the whole election key (see Section~\ref{sec:vocdoni-protocol:flow:keys}).

\paragraph{Voters.} Voters are the participants that belong to the census and are allowed to cast their votes in the election. Each voter holds a secp256k1 key pair, whose Ethereum address identifies the voter in the census, and is assigned a weight, which bounds or scales their vote according to the ballot rules. Voters may cast as many ballots as they wish during the voting period; each new ballot overwrites the previous one.

\paragraph{Sequencers.} Sequencers are a set of parties that during the voting period receive and verify encrypted ballots from voters, reencrypt votes, and update the shared public state accordingly. A sequencer groups valid ballots into batches, proves each batch with the batch program in a \ac{zkVM}, publishes the batch data in blobs, and submits the proof to the process registry. The role is permissionless: any party with an Ethereum account can run a sequencer for any election, and several sequencers can serve the same election concurrently, in which case they race to settle the next state transition. Proof generation is computationally intensive and is typically delegated to a GPU \emph{prover} service. The prover learns the private inputs of the batch, including the reencryption secrets, so a sequencer must run its own prover or trust its operator with ballot privacy (but not with integrity, which the proof guarantees).

\paragraph{Observers.} Any party can run a sequencer node without a signing key. Such a node, called an observer, never settles transitions, but it replays every transition settled by others from the published blobs, checks that the resulting state roots match the ones recorded on-chain, and serves the state of the election, including the inclusion proofs that voters use to check that their ballots were recorded. Observers therefore provide independent verification of the election and of its tally.
